\documentclass[11pt,letterpaper]{article}

\usepackage[margin=0.85in]{geometry}
\usepackage{fontspec}
\usepackage{xeCJK}
\setmainfont{Times New Roman}
\setCJKmainfont{Microsoft YaHei}
\usepackage{hyperref}
\usepackage{enumitem}
\usepackage{setspace}
\usepackage{tabularx}
\usepackage{array}

\hypersetup{colorlinks=true, allcolors=black}
\setlist[itemize]{leftmargin=1.2em, itemsep=0.25em, topsep=0.15em, parsep=0.05em}
\setlength{\parindent}{0pt}
\setlength{\parskip}{0.15em}
\pagestyle{empty}
\setstretch{1.15}

\newlength{\cvdatecol}
\setlength{\cvdatecol}{11em}

\newcommand{\sectionheader}[1]{%
  \vspace{0.45em}
  {\large\bfseries #1}\\[-0.15em]
  \rule{\textwidth}{0.4pt}\par
  \vspace{0.2em}
}

\newcommand{\cvline}[3][0em]{%
  \par\nobreak
  \noindent\hspace*{#1}%
  \begin{tabularx}{\dimexpr\textwidth-#1\relax}{@{}X@{\hspace{0.75em}}>{\raggedleft\arraybackslash}p{\cvdatecol}@{}}
  #2 & #3 \\
  \end{tabularx}%
}

\newcommand{\courseline}[1]{%
  \par\addvspace{0.22em}%
  \noindent\hspace*{3.2em}#1\par
}

\begin{document}

\begin{center}
{\LARGE\bfseries 許可證（Ke Cheng HSU, Clio）}\\[0.5em]
\href{mailto:khsu@connect.ust.hk}{khsu@connect.ust.hk}\\[0.15em]
\href{https://cliohsu1997.github.io}{https://cliohsu1997.github.io/}
\end{center}

\vspace{0.45em}

\textbf{通讯地址}\\[0.15em]
\hspace{1.5em} 香港九龙清水湾香港科技大学李兆基商学大楼 5067 室\\
\hspace{1.5em} 经济系

\vspace{0.4em}

\textbf{研究方向}\\[0.15em]
\hspace{1.5em} 产业组织、数字经济、量化营销

\sectionheader{教育背景}
{\setstretch{0.88}
香港科技大学\\[-0.3em]
\cvline[1.5em]{经济学博士}{2022 -- 2027（预计）}
\vspace{-0.12em}
\cvline[1.5em]{经济学哲学硕士（M.Phil.）}{2020 -- 2022}
\vspace{-0.12em}
\cvline[1.5em]{数学与经济学理学学士（副修精算数学）}{2015 -- 2019}
}

\vspace{-0.4em}
\sectionheader{工作论文}
\begin{itemize}
  \item 《电商退货政策设计：购买效率与退货效率的权衡》（Designing Return Policy in E-Commerce: A Trade-off between Purchase and Return Efficiency）\textit{（求职论文）}\\[0.15em]
  \textit{摘要：}本文研究退货政策是否同时促进有效退货与有效购买。更严格的退货政策可以抑制事后无效退货，但退货率下降会使卖家更倾向于提高价格，从而可能把事前有效的购买排除在外。我们在六个服装品类、垄断定价下估计购买、运费险与退货决策的结构模型，比较自动报销退货运费的保障（RI）与可选保障（CI）。宽松政策（如 RI）在高毛利、高残值品类中表现更好，因为它扩大了有价值的购买；严格政策（如 CI）在毛利较薄的品类中表现更好，因为它抑制过多退货。部分退货罚金（Partial RPD）——将折扣价与中等程度补货费配对的菜单——在多数品类中比 RI 和 CI 更好地平衡购买效率与退货效率。\\[0.15em]
  \textit{报告（$^{\ast}$ 为已安排）：} APIOC 2026$^{\ast}$；2026 Econometric Society Australasia Meeting（ESAM 2026）$^{\ast}$；中国经济学年会夏季论坛 2026；香港科技大学商学院博士生会议；香港科技大学 CEP Brownbag。

  \item 《异质搜寻成本下的产品设计》（Product Design with Heterogeneous Search Costs），与 Pak Hung Au\\[0.15em]
  \textit{摘要：}在 Wolinsky（1986）差异化产品序贯搜寻模型的基础上，本文引入搜寻成本异质性，分析企业如何利用可观察的产品设计影响消费者的搜寻行为。均衡中，消费者的搜寻呈现分割结构：搜寻成本较低者倾向于寻找利基产品，搜寻成本较高者倾向于标准化产品。在两种产品设计下，若搜寻模式的内生转变（构成效应）主导了鼓励搜寻的效应，搜寻成本下降可能使两种产品的价格同时上升。此外，价格对搜寻成本下降的反应取决于冲击落在成本分布中的位置，因而并不对称。\\[0.15em]
  \textit{报告：} APIOC 2023；AMES 2023，新加坡南洋理工大学。
\end{itemize}

\sectionheader{进行中的研究}
\begin{itemize}
  \item 《含退货的价格转嫁与税负归宿》（Pass-Through and Incidence with Product Returns），与 Ce Bian、Liang Zhao\\[0.15em]
  \textit{摘要：}本文研究存在产品退货时的价格转嫁与税负归宿。与无退货情形不同，平台佣金或销售税可以产生负向转嫁并提高社会福利。当高价格显著提高消费者的退货概率时，有效利润率下降，卖家减少供给。这条向下倾斜的供给曲线可以使均衡价格随税收上升而下降，并通过降低退货率提高社会剩余。
  \item 《滞销过多？价格承诺的环境后果》（Too Much Unsold Stock?: Environmental Consequences of Price Commitment），与 Kohei Kawaguchi、Kei Kawai\\[0.15em]
  \textit{摘要：}转售价格维持具有促进竞争的作用，因为它鼓励卖家在低需求状态下持有库存。但当实现需求低于库存时，卖家不能降价出清，未售出的商品被退回并带来环境成本。我们构建含需求不确定性的动态寡头模型，用月度书店数据估计，并考察有多少滞销来自价格承诺，以及押金—退款税如何减少浪费。
\end{itemize}

\sectionheader{经历}
{\setstretch{0.88}
\textit{研究}\\[-0.4em]
\cvline[1.5em]{研究助理，Benjamin Bernard 教授，台湾大学}{2020}

\vspace{-0.05em}
\textit{教学}\\[-0.4em]
\cvline[1.5em]{讲师，经济学硕士数学营，香港科技大学（4.53/5）}{2026}
\vspace{-0.1em}
\cvline[1.5em]{教学助理协调员，商学院，香港科技大学}{2025 -- 2026}
\vspace{-0.1em}
\cvline[1.5em]{教学助理，经济系，香港科技大学（平均 4.72/5）}{2020 -- 2025}
\vspace{0.05em}
\courseline{ECON2123 宏观经济学（4.53/5）}
\courseline{ECON4134 不确定性与信息经济学（4.83/5）}
\courseline{ECON4999T 不确定性与信息经济学（4.56/5）}
\courseline{ECON5210 微观经济理论 I（4.85/5；4.82/5）}
}

\sectionheader{奖助}
{\setstretch{0.9}
\cvline[1.5em]{研究差旅资助}{2023；2026}
\vspace{-0.1em}
\cvline[1.5em]{Red Bird 博士奖}{2022}
\vspace{-0.1em}
\cvline[1.5em]{香港科技大学研究生助学金}{2020 -- 至今}
\vspace{-0.1em}
\cvline[1.5em]{学术成就奖章}{2019}
\vspace{-0.1em}
\cvline[1.5em]{大学奖学金（73 万港元）}{2015 -- 2019}
}

\sectionheader{技能}
技术：R、Python、C++、\LaTeX、MATLAB、AWS\\
语言：英语（流利）、普通话（母语）、粤语（基础）

\sectionheader{推荐人（论文指导委员会）}
{\setstretch{0.92}
\noindent
\begin{minipage}[t]{0.48\textwidth}
\textbf{\href{https://aupakhung.weebly.com/}{Pak Hung Au}}\\
副教授\\
经济系\\
香港科技大学\\
\href{mailto:aupakhung@ust.hk}{aupakhung@ust.hk}
\end{minipage}\hfill
\begin{minipage}[t]{0.48\textwidth}
\textbf{\href{https://www.kohei-kawaguchi.com/}{Kohei Kawaguchi}}\\
副教授\\
经济系\\
香港科技大学\\
\href{mailto:kkawaguchi@ust.hk}{kkawaguchi@ust.hk}
\end{minipage}

\vspace{0.85em}
\noindent
\begin{minipage}[t]{0.48\textwidth}
\textbf{\href{https://www.e.u-tokyo.ac.jp/fservice/faculty/kawai/kawai-e/kawai01-e.html}{Kei Kawai}}\\
教授\\
经济系\\
东京大学\\
\href{mailto:kawaixkei@gmail.com}{kawaixkei@gmail.com}
\end{minipage}
}

\end{document}
